\documentclass[11pt,letterpaper]{article}
\usepackage{fontspec}
\usepackage{tocloft}
\usepackage{multicol}
\usepackage{fancyhdr}
\usepackage{nicefrac}
\usepackage[left=1.3in,right=1.3in,top=1in,bottom=1in]{geometry}
\setmainfont[Scale=1.2,AutoFakeBold=1.5,AutoFakeSlant=0.3]{Doves Type}

\title{Pork Baby Back Ribs}
\author{}
\date{}

\pagestyle{fancy}
\fancyhf{}
\fancyhead[C]{\textit{Pork Baby Back Ribs}}
\fancyfoot[C]{\thepage}
\renewcommand{\headrulewidth}{0pt}
\renewcommand{\footrulewidth}{0pt}

\begin{document}

\maketitle
\thispagestyle{empty}

\section*{Ingredients}
\setlength{\columnsep}{20pt}
\begin{multicols}{2}
\noindent
    Baby back ribs \dotfill 3\nicefrac{1}{2}~lb. (1 rack) \\
    Kosher salt, coarse \dotfill 2 tsp. \\
    Black pepper, coarse \dotfill 4 tsp. \\
    Smoked paprika \dotfill 1 tsp. \\
    Garlic powder \dotfill \nicefrac{1}{2} tsp. \\
    Smoked chile powder \dotfill \nicefrac{1}{4} tsp. \\
    \columnbreak
    Yellow onion, halved \dotfill 1 medium \\
    Garlic cloves, smashed \dotfill 4 \\
    Bay leaves \dotfill 2 \\
    Water \dotfill 1 cup \\
    Apple cider vinegar \dotfill \nicefrac{1}{2} cup \\
    Hickory smoke concentrate \dotfill 1 tsp. \\
    BBQ sauce (see p.~\pageref{sec:sauce}) \dotfill \nicefrac{1}{2}--\nicefrac{3}{4} cup \\
\end{multicols}

\section*{Directions}

\noindent
Combine 2~tsp. \textbf{kosher salt}, 4~tsp. \textbf{black pepper}, 1~tsp. \textbf{smoked paprika}, \nicefrac{1}{2}~tsp. \textbf{garlic powder}, and \nicefrac{1}{4}~tsp. \textbf{smoked chile powder} in \textit{Small Bowl~\#1} (rub) ---
Remove membrane from bone side of \textbf{ribs}: slide a butter knife under the thin silverskin at one end, grip with a paper towel, and pull the length of the rack in one motion ---
Pat \textbf{ribs} dry and apply rub (\textit{Small Bowl~\#1}) liberally to all surfaces; rest \textit{30~minutes} at room temperature, or refrigerate uncovered up to \textit{24~hours} ---
Halve \textbf{onion} and smash \textbf{garlic}; combine with \textbf{bay leaves} in \textit{Small Bowl~\#2} (aromatics) ---
Combine 1~cup \textbf{water}, \nicefrac{1}{2}~cup \textbf{apple cider vinegar}, and 1~tsp. \textbf{hickory smoke} in \textit{Medium Bowl~\#1} (braising liquid)

\begin{enumerate}
    \item Pour braising liquid (\textit{Medium Bowl~\#1}) into the Instant Pot. Add aromatics (\textit{Small Bowl~\#2}). Place the trivet in the pot.

    \item Cut the \textbf{rib rack} in half crosswise. Stand the two sections on the trivet, bone side down, curved against the interior wall of the pot.

    \item Seal the lid and set the venting knob to Sealing. Pressure cook on High for \textit{25~minutes}. Allow a \textit{10-minute} natural pressure release, then quick-release any remaining pressure.

    \item Carefully transfer \textbf{ribs} to a foil-lined tray or cutting board---they will be fragile. Strain all \textbf{braising liquid} into \textit{Large Bowl~\#1} and reserve for the sauce. \textbf{Make the BBQ sauce now} (p.~\pageref{sec:sauce}) before proceeding---the air fryer step requires it.

    \item Brush \textbf{ribs} on all exposed surfaces with a generous coat of \textbf{BBQ sauce}. Air fry bone side up for \textit{4~minutes} until the glaze begins to char and blister.

    \item Flip \textbf{ribs} to meaty side up. Brush with a second coat of \textbf{BBQ sauce}. Air fry for \textit{4--5~minutes} until the glaze is deeply caramelized and blackened at the edges.

    \item Rest \textit{5~minutes}, then slice between bones and serve with remaining \textbf{BBQ sauce} on the side.
\end{enumerate}

\newpage

{\footnotesize
\setlength{\columnsep}{20pt}
\setlength{\multicolsep}{6pt}
\begin{multicols}{2}
\setlength{\parindent}{0pt}
\setlength{\parskip}{2pt}
\setlength{\itemsep}{0pt}
\setlength{\parsep}{0pt}

\subsection*{Equipment Required}
\begin{itemize}
    \item 6-quart or larger Instant Pot
    \item Instant Pot trivet or steam rack
    \item Air fryer (basket-style preferred)
    \item Foil-lined sheet tray or cutting board
    \item Basting brush
    \item Fat separator or ladle (for sauce)
    \item Medium saucepan (for sauce)
    \item Small prep bowls (2)
    \item Medium prep bowl (1)
    \item Large bowl (1, for braising liquid)
    \item Butter knife and paper towels (membrane removal)
    \item Tongs (2 pairs recommended)
    \item Measuring cups and spoons
\end{itemize}

\subsection*{Mise en Place}
\begin{itemize}
    \item Small Bowl \#1 --- rub: 2~tsp. \textbf{kosher salt}, 4~tsp. \textbf{black pepper}, 1~tsp. \textbf{smoked paprika}, \nicefrac{1}{2}~tsp. \textbf{garlic powder}, \nicefrac{1}{4}~tsp. \textbf{smoked chile powder}
    \item Small Bowl \#2 --- aromatics: halved \textbf{onion}, smashed \textbf{garlic}, \textbf{bay leaves}
    \item Medium Bowl \#1 --- braising liquid: 1~cup \textbf{water}, \nicefrac{1}{2}~cup \textbf{apple cider vinegar}, 1~tsp. \textbf{hickory smoke concentrate}
    \item Large Bowl \#1 --- reserved \textbf{braising liquid} after cooking (for sauce)
    \item Apply rub and rest \textbf{ribs} \textit{30~minutes} before cooking, or overnight in the refrigerator
\end{itemize}

\subsection*{Ingredient Tips}
\begin{itemize}
    \item Choose \textbf{ribs} with a full, even layer of meat across the rack; avoid thin or pre-cut racks
    \item Use coarse-grind \textbf{black pepper} (16-mesh or equivalent)---fine-grind turns acrid under pressure; coarse grind builds bark character
    \item \textbf{Kosher salt} (Diamond Crystal preferred) distributes more evenly than table salt at this volume
    \item \textbf{Hickory smoke concentrate} (e.g., Colgin, Wright's) is more potent than standard liquid smoke; 1~tsp. is sufficient---do not overdo it
    \item \textbf{Apple cider vinegar} brightens the braising liquid and tenderizes the meat; do not substitute white vinegar
\end{itemize}

\subsection*{Preparation Tips}
\begin{itemize}
    \item Membrane removal is non-negotiable---leaving it on creates a rubbery barrier that blocks rub penetration and produces an unpleasant texture
    \item Overnight dry brine (rub applied, refrigerated uncovered) produces noticeably deeper flavor and better surface texture in the air fryer
    \item Standing the \textbf{ribs} on the trivet rather than coiling them on the bottom ensures even pressure cooking and prevents the bottom from overcooking in the liquid
    \item The ribs will be very tender after pressure cooking; use two sets of tongs or a wide spatula to transfer them without the rack splitting
    \item Do not skip the air fryer rest between sides---pulling the ribs too quickly risks them breaking apart
    \item Two sauce coats are required: the first coat carbonizes into the meat; the second coat provides the lacquered glaze finish
\end{itemize}

\subsection*{Make Ahead \& Storage}
\begin{itemize}
    \item Apply rub up to \textit{24~hours} ahead and refrigerate uncovered
    \item Pressure cook up to \textit{4~hours} ahead; hold \textbf{ribs} loosely covered at room temperature and air fry to order
    \item Do not sauce \textbf{ribs} until just before the air fryer step---sauce applied too early will steam rather than caramelize
    \item Leftovers keep \textit{3--4~days} refrigerated; reheat in air fryer at \textit{375°F} for \textit{5--7~minutes}
    \item Freeze unsauced pressure-cooked \textbf{ribs} up to \textit{2~months}; thaw overnight and finish in air fryer from cold
\end{itemize}

\subsection*{Serving Suggestions}
\begin{itemize}
    \item Serve immediately off the air fryer with extra \textbf{BBQ sauce} on the side
    \item Classic accompaniments: Texas toast, pickled jalape\~nos, white onion, dill pickles
    \item Pairs well with pinto beans, coleslaw, or elote
    \item For a platter presentation, slice between every bone and arrange meat side up; pool sauce underneath
\end{itemize}

\end{multicols}
}
\end{document}